\bookchapter{付録C：ミニマル用語「隠語」対照表}

\begin{center}
\begin{tabularx}{\linewidth}{|X|X|X|}
\hline
\textbf{：問題に遭遇したら、まず再起動！これは半分以上のスマートフォンの小さな不具合を解決する万能キーです。} & \textbf{付録C：ミニマル用語「隠語」対照表} & \textbf{聞いたことがあるかもしれない言葉} \tabularnewline
\hline
実はこれのこと…… & 一言で説明 & SoC / チップ \tabularnewline
スマートフォンの脳と心臓 & スマートフォンの速さやゲームのカクつきを決めるコアパーツ。 & RAM（メモリ） \tabularnewline
スマートフォンの「短期作業デスク」 & デスクが大きいほど、同時に開けるアプリが多く、切り替えもスムーズ。 & ROM（ストレージ） \tabularnewline
スマートフォンの「大倉庫」 & 写真、動画、アプリを保存する場所。128G、256Gと言っているのはこれのこと。 & OLED / LCD \tabularnewline
画面の2種類の素材 & OLEDは色鮮やかで省電力、LCDは目に優しくコストが低い。 & NFC \tabularnewline
近距離無線通信 & スマートフォンを交通系ICカード、入退室カード、銀行カードに変える「タッチ」機能。 & OTAアップデート \tabularnewline
ワイヤレスシステムアップデート & 【設定】でタップするあの「システムアップデート」がこれです。 & Root \tabularnewline
Androidスーパー管理者権限 & 取得するとシステムファイルを変更できますが、リスクもあります。 & BLロック \tabularnewline
Bootloaderロック & メーカーがかけるセキュリティロック、解除して初めてROM焼きが可能。 & NPU \tabularnewline
ニューラルネットワークプロセッサ & AIタスクに特化したチップ、CPUより数十倍速く、消費電力も低い。 & オンデバイスAI \tabularnewline
スマートフォン上でローカル実行されるAI & AIがスマートフォン上で直接実行され、ネット接続不要、高速でプライバシーに優れる。 & Qi2 \tabularnewline
次世代ワイヤレス充電規格 & マグネット位置決め付きのワイヤレス充電、MagSafeのようなもので、充電がより便利に。 & eSIM \tabularnewline
組み込み型SIMカード & 実体カードのないSIM、スマートフォンに直接書き込み、キャリア変更時にカードの差し替え不要。 & LTPO \tabularnewline
\hline
\end{tabularx}
\end{center}
